\documentclass[10pt,a4paper]{amsart}
\usepackage{amsmath,amssymb,amsthm,graphicx}
\usepackage{verbatim}
\newenvironment{reference}{\comment}{\endcomment}
%\newenvironment{reference}{}{}
\newenvironment{slogan}{\comment}{\endcomment}
\newenvironment{history}{\comment}{\endcomment}

% For commutative diagrams we use Xy-pic
\usepackage[all]{xy}

% We use 2cell for 2-commutative diagrams.
\xyoption{2cell}
% Theorem environments.
%
\theoremstyle{plain}
\newtheorem{theorem}[subsection]{Theorem}
\newtheorem{proposition}[subsection]{Proposition}
\newtheorem{lemma}[subsection]{Lemma}

\theoremstyle{definition}
\newtheorem{definition}[subsection]{Definition}
\newtheorem{example}[subsection]{Example}
\newtheorem{exercise}[subsection]{Exercise}
\newtheorem{situation}[subsection]{Situation}

\theoremstyle{remark}
\newtheorem{remark}[subsection]{Remark}
\newtheorem{remarks}[subsection]{Remarks}

\numberwithin{equation}{subsection}

% Macros
%
\def\lim{\mathop{\mathrm{lim}}\nolimits}
\def\colim{\mathop{\mathrm{colim}}\nolimits}
\def\Spec{\mathop{\mathrm{Spec}}}
\def\Hom{\mathop{\mathrm{Hom}}\nolimits}
\def\Ext{\mathop{\mathrm{Ext}}\nolimits}
\def\SheafHom{\mathop{\mathcal{H}\!\mathit{om}}\nolimits}
\def\SheafExt{\mathop{\mathcal{E}\!\mathit{xt}}\nolimits}
\def\Sch{\mathit{Sch}}
\def\Mor{\mathop{\mathrm{Mor}}\nolimits}
\def\Ob{\mathop{\mathrm{Ob}}\nolimits}
\def\Sh{\mathop{\mathit{Sh}}\nolimits}
\def\NL{\mathop{N\!L}\nolimits}
\def\CH{\mathop{\mathrm{CH}}\nolimits}
\def\proetale{{pro\text{-}\acute{e}tale}}
\def\etale{{\acute{e}tale}}
\def\QCoh{\mathit{QCoh}}
\def\Ker{\mathop{\mathrm{Ker}}}
\def\Im{\mathop{\mathrm{Im}}}
\def\Coker{\mathop{\mathrm{Coker}}}
\def\Coim{\mathop{\mathrm{Coim}}}

% Boxtimes
%
\DeclareMathSymbol{\boxtimes}{\mathbin}{AMSa}{"02}

%
% Macros for moduli stacks/spaces
%
\def\QCohstack{\mathcal{QC}\!\mathit{oh}}
\def\Cohstack{\mathcal{C}\!\mathit{oh}}
\def\Spacesstack{\mathcal{S}\!\mathit{paces}}
\def\Quotfunctor{\mathrm{Quot}}
\def\Hilbfunctor{\mathrm{Hilb}}
\def\Curvesstack{\mathcal{C}\!\mathit{urves}}
\def\Polarizedstack{\mathcal{P}\!\mathit{olarized}}
\def\Complexesstack{\mathcal{C}\!\mathit{omplexes}}
% \Pic is the operator that assigns to X its picard group, usage \Pic(X)
% \Picardstack_{X/B} denotes the Picard stack of X over B
% \Picardfunctor_{X/B} denotes the Picard functor of X over B
\def\Pic{\mathop{\mathrm{Pic}}\nolimits}
\def\Picardstack{\mathcal{P}\!\mathit{ic}}
\def\Picardfunctor{\mathrm{Pic}}
\def\Deformationcategory{\mathcal{D}\!\mathit{ef}}
\usepackage{cleveref}
\crefname{originalreference}{original reference}{original references}
\Crefname{originalreference}{Original reference}{Original references}
\makeatletter
\newlabel{definition-gysin-homomorphism}{{02T8}{source}{}{sourcepointer-1}{}}
\newlabel{definition-gysin-homomorphism@cref}{{[originalreference][0][]02T8}{[1][0][]source}{}{}{}}
\makeatother
\crefname{originalreference}{original reference}{original references}
\Crefname{originalreference}{Original reference}{Original references}
\makeatletter
\newlabel{lemma-closed-in-X-gysin}{{02TA}{source}{}{sourcepointer-2}{}}
\newlabel{lemma-closed-in-X-gysin@cref}{{[originalreference][0][]02TA}{[1][0][]source}{}{}{}}
\makeatother
\crefname{originalreference}{original reference}{original references}
\Crefname{originalreference}{Original reference}{Original references}
\makeatletter
\newlabel{lemma-factors}{{02TI}{source}{}{sourcepointer-3}{}}
\newlabel{lemma-factors@cref}{{[originalreference][0][]02TI}{[1][0][]source}{}{}{}}
\makeatother
\crefname{originalreference}{original reference}{original references}
\Crefname{originalreference}{Original reference}{Original references}
\makeatletter
\newlabel{lemma-gysin-factors-general}{{02TM}{source}{}{sourcepointer-4}{}}
\newlabel{lemma-gysin-factors-general@cref}{{[originalreference][0][]02TM}{[1][0][]source}{}{}{}}
\makeatother
\crefname{originalreference}{original reference}{original references}
\Crefname{originalreference}{Original reference}{Original references}
\makeatletter
\newlabel{lemma-proper-pushforward-rational-equivalence}{{02S2}{source}{}{sourcepointer-5}{}}
\newlabel{lemma-proper-pushforward-rational-equivalence@cref}{{[originalreference][0][]02S2}{[1][0][]source}{}{}{}}
\makeatother
\crefname{originalreference}{original reference}{original references}
\Crefname{originalreference}{Original reference}{Original references}
\makeatletter
\newlabel{remark-pullback-pairs}{{0B6Y}{source}{}{sourcepointer-6}{}}
\newlabel{remark-pullback-pairs@cref}{{[originalreference][0][]0B6Y}{[1][0][]source}{}{}{}}
\makeatother
\crefname{originalreference}{original reference}{original references}
\Crefname{originalreference}{Original reference}{Original references}
\makeatletter
\newlabel{situation-setup}{{02QL}{source}{}{sourcepointer-7}{}}
\newlabel{situation-setup@cref}{{[originalreference][0][]02QL}{[1][0][]source}{}{}{}}
\makeatother
\title{The Gysin homomorphism factors through rational equivalence}
\author{The Stacks Project Authors}
\date{stacks-project:spaces-chow.tex:local-snapshot-20261012; GNU FDL 1.2 or later}
\begin{document}
\maketitle
\noindent\textit{Editorial note.} All delivered mathematical text is the author's own, copied from the pinned Stacks Project chapter file spaces-chow.tex; no statement or proof step was rewritten, reordered or removed.\par\medskip
\noindent\textit{Editorial note.} The counted proof is the author's complete argument: a coproduct and proper pushforward reduction to a single rational function on an integral space, followed by a two case identification with Chern class caps\par\medskip
\noindent\textit{Editorial note.} The label lemma-gysin-factors names the counted statement, and the proof environment immediately following it in the source is delivered verbatim as its proof.\par\medskip
\noindent\textit{Editorial note.} The counted statement is carried by the author's own lemma environment in the chapter source.\par\medskip
\noindent\textit{Allowed standard prerequisites (editorial).}\begin{itemize}
\item The Gysin homomorphism attached to a triple consisting of an invertible sheaf, a nonzero meromorphic section and its effective Cartier divisor of zeros, and the compatibility of Gysin pullback with proper pushforward
\item Proper pushforward of cycles preserves rational equivalence, so that a vanishing result on a proper cover descends
\item The general Gysin and first Chern class compatibility, and the vanishing of the cap product with the first Chern class of the structure sheaf
\end{itemize}\medskip
\ifdefined\hypertarget\hypertarget{sourcepointer-1}{}\fi
\noindent\textit{Source reference.} Original reference 02T8: sources/stacks-spaces-chow/pinned-original-source.tex, source lines 3448--3448. Prerequisite pointer invoked at proof line(s) 3733; it names a standard fact of the same chapter that the counted argument uses, and every step of the counted proof built on it is retained verbatim in the delivered body.\par\medskip
\ifdefined\hypertarget\hypertarget{sourcepointer-2}{}\fi
\noindent\textit{Source reference.} Original reference 02TA: sources/stacks-spaces-chow/pinned-original-source.tex, source lines 3545--3545. Prerequisite pointer invoked at proof line(s) 3752; it names a standard fact of the same chapter that the counted argument uses, and every step of the counted proof built on it is retained verbatim in the delivered body.\par\medskip
\ifdefined\hypertarget\hypertarget{sourcepointer-3}{}\fi
\noindent\textit{Source reference.} Original reference 02TI: sources/stacks-spaces-chow/pinned-original-source.tex, source lines 3327--3327. Prerequisite pointer invoked at proof line(s) 3763; it names a standard fact of the same chapter that the counted argument uses, and every step of the counted proof built on it is retained verbatim in the delivered body.\par\medskip
\ifdefined\hypertarget\hypertarget{sourcepointer-4}{}\fi
\noindent\textit{Source reference.} Original reference 02TM: sources/stacks-spaces-chow/pinned-original-source.tex, source lines 3662--3662. Prerequisite pointer invoked at proof line(s) 3766; it names a standard fact of the same chapter that the counted argument uses, and every step of the counted proof built on it is retained verbatim in the delivered body.\par\medskip
\ifdefined\hypertarget\hypertarget{sourcepointer-5}{}\fi
\noindent\textit{Source reference.} Original reference 02S2: sources/stacks-spaces-chow/pinned-original-source.tex, source lines 2176--2176. Prerequisite pointer invoked at proof line(s) 3754; it names a standard fact of the same chapter that the counted argument uses, and every step of the counted proof built on it is retained verbatim in the delivered body.\par\medskip
\ifdefined\hypertarget\hypertarget{sourcepointer-6}{}\fi
\noindent\textit{Source reference.} Original reference 0B6Y: sources/stacks-spaces-chow/pinned-original-source.tex, source lines 3496--3496. Prerequisite pointer invoked at proof line(s) 3751; it names a standard fact of the same chapter that the counted argument uses, and every step of the counted proof built on it is retained verbatim in the delivered body.\par\medskip
\ifdefined\hypertarget\hypertarget{sourcepointer-7}{}\fi
\noindent\textit{Source reference.} Original reference 02QL: sources/stacks-spaces-chow/pinned-original-source.tex, source lines 264--264. Prerequisite pointer invoked at proof line(s) 3731; it names a standard fact of the same chapter that the counted argument uses, and every step of the counted proof built on it is retained verbatim in the delivered body.\par\medskip
\begin{lemma}
\label{lemma-gysin-factors}
In Situation \ref{situation-setup} let $X/B$ be good.
Let $(\mathcal{L}, s, i : D \to X)$ be as in
Definition \ref{definition-gysin-homomorphism}.
The Gysin homomorphism factors through rational equivalence to
give a map $i^* : \CH_{k + 1}(X) \to \CH_k(D)$.
\end{lemma}

\begin{proof}
Let $\alpha \in Z_{k + 1}(X)$ and assume that $\alpha \sim_{rat} 0$.
This means there exists a locally finite collection of integral
closed subspaces $W_j \subset X$ of $\delta$-dimension $k + 2$
and $f_j \in R(W_j)^*$ such that
$\alpha = \sum i_{j, *}\text{div}_{W_j}(f_j)$.
Set $X' = \coprod W_i$ and consider the diagram
$$
\xymatrix{
D' \ar[d]_q \ar[r]_{i'} & X' \ar[d]^p \\
D \ar[r]^i & X
}
$$
of Remark \ref{remark-pullback-pairs}. Since $X' \to X$ is proper
we see that $i^*p_* = q_*(i')^*$ by Lemma \ref{lemma-closed-in-X-gysin}.
As we know that $q_*$ factors through rational equivalence
(Lemma \ref{lemma-proper-pushforward-rational-equivalence}), it suffices
to prove the result for $\alpha' = \sum \text{div}_{W_j}(f_j)$
on $X'$. Clearly this reduces us to the case where $X$ is integral
and $\alpha = \text{div}(f)$ for some $f \in R(X)^*$.

\medskip\noindent
Assume $X$ is integral and $\alpha = \text{div}(f)$ for some $f \in R(X)^*$.
If $X = D$, then we see that $i^*\alpha$ is equal
to $c_1(\mathcal{L}) \cap \alpha$.
This is rationally equivalent to zero by Lemma \ref{lemma-factors}.
If $D \not = X$, then we see that $i^*\text{div}_X(f)$ is equal to
$c_1(\mathcal{O}_D) \cap [D]_{n - 1}$ in $\CH_k(D)$ by
Lemma \ref{lemma-gysin-factors-general}. Of course
capping with $c_1(\mathcal{O}_D)$ is the zero map.
\end{proof}

\end{document}
